\documentclass[11pt,a4paper]{article}
\usepackage[margin=23mm]{geometry}
\usepackage{amsmath,amssymb,amsthm}
\setlength{\emergencystretch}{2em}
\newtheorem{lemma}{Lemma}
\newtheorem{theorem}[lemma]{Theorem}
\newcommand{\gen}[1]{\langle#1\rangle}
\title{A transitive witness and its transported orbit}
\author{Human Compression project: L044 continuation}
\date{}
\begin{document}\maketitle
Let $G\le\operatorname{Alt}(\Omega)$ be transitive with
connected string C-marking $\rho_0,\ldots,\rho_{r-1}$,
all maximal omissions intransitive, $r\ge6$ and $n=2r+1$.
The first three marks are supported on, and transitive on,
the actual six-point head $D$. Retain the literal interfaces
\[
 D=\{1,\ldots,6\},\quad \tau=\rho_3=(6\ \beta)\delta,
 \quad \sigma=\rho_4,\quad H=G_{\ge5},\quad K=G_{>5}.
\]
Here $G_{\ge3}=\operatorname{Alt}(\Omega\setminus\{1,\ldots,5\})$,
$\delta$ is an odd involution supported outside $D\cup\{\beta\}$,
and $\sigma$ fixes $D$. The group $H$ fixes the seven head points
$D\cup\{\beta\}$. Every colour $j\ge5$ has two edges joining
distinct orbits of its full ambient omission $G_j$; these
are called ambient eligible. The actual anchor
$z_0=\beta^\sigma$ lies outside those seven points,
is fixed by $K$ and by $\delta$, and is not $H$-fixed.
The degree-independent natural alternating overlap proof
also gives a pure three-cycle on $\{6,\beta,z_0\}$
in $\gen{\tau,\sigma}$.

Choose $j$ by the complete strong simultaneous-witness
theorem. Each actual $H$-orbit has at most one
ambient-eligible $j$-edge. Every witness orbit on which
$K$ is intransitive has a full local unique crossing
and is fixed pointwise by $\delta$.
Suppose some witness orbit $O$ is transitive under $K$.
Necessarily $j\ge6$, since an ambient $5$-edge cannot
have both endpoints in a transitive $K$-orbit.
The complete whole-orbit transport supplies
$p\in G_{<5}$ for which
\[
                         T=O^p
\]
is an actual $K$-orbit in a $K$-intransitive $H$-orbit
$W$. The transport commutes with $K$, preserves the
ambient witness edge, and transports the complete marked
$K$-action. In particular $O,T$ are disjoint and
$\delta$ fixes $T$ pointwise.

\begin{lemma}[The transported action has enough points]
The active $K$-labels on $O$ and $T$ form an interval
$6,\ldots,x$ for some $x\ge j$. Every omission of a
label in $6,\ldots,x-1$ is intransitive in this actual
restricted action. Consequently
\[
                         |O|=|T|\ge x-4.
\]
No intersection property for the restricted action is assumed.
\end{lemma}
\begin{proof}
First apply the complete commuting-component argument
at the unique ambient $j$-edge to the transitive
$K$-action on $O$. For clarity, any component of its
noncommutation diagram not containing $j$ commutes
with $\rho_j$, preserves $O$ and lies in $G_j$.
It must fix both endpoints of the unique eligible
edge. A factor of a transitive commuting product
which fixes a point is trivial everywhere. Thus
only the component containing $j$ can be active.
The active labels form a connected interval $m,\ldots,x$,
with $m\ge6$.

If $m\ge7$, then $\rho_5$ centralizes the whole
$K$-action on $O$. It preserves this actual $H$-orbit
and fixes both endpoints of its unique ambient
$j$-edge, since $j\ge7$. It therefore fixes all
of $O$. The active $H$-interval on this witness
would start at $m\ge7$, contrary to the complete
acquired first-active-label theorem. Hence $m=6$.
Notice that this argument does not assume
$\rho_5$ fixes a general $K$-transitive witness.

The omission of $6$ is intransitive. If $j=6$,
its two ambient witness endpoints already lie in
distinct orbits of that omission. If $j\ge7$,
work on the transported orbit $T\subset W$.
The generator $\rho_5$ fixes the transported edge
endpoints, because $W$ has its unique ambient
$j$-edge and $\rho_5\in G_j$ commutes with
$\rho_j$. If $G_{\ge7}$ were transitive on $T$,
its centralization by $\rho_5$ would force
$\rho_5$ to fix all of $T$. Then both $K$
and $\rho_5$ would preserve $T$, making $W=T$,
contrary to the $K$-intransitivity of $W$.

For an internal label $6<i<x$, its omission is
also intransitive. For $i=j$ this follows from
ambient eligibility. Otherwise, suppose the
omission were transitive. Its left and right
factors commute and are both nontrivial, since
all labels $6,\ldots,x$ are active. The factor
on the side not containing $j$ centralizes
$\rho_j$ and fixes both endpoints of the unique
ambient $j$-edge in $O$. Transitivity of the
commuting product would make that factor trivial,
a contradiction. This proves every asserted omission.

Put $d=x-5$, the number of active marks. For each
of the first $d-1$ labels choose an edge joining
distinct orbits of its full restricted omission.
These selected edges form a forest: a cycle would
join the endpoints of any one of its edges using
only the other labels, contradicting eligibility.
Thus the degree is at least $d$. If it were exactly
$d$, the forest would be a spanning tree. Deleting
one selected edge divides the tree into two connected
sets. Its intransitive full omission contains every
other selected edge, so its orbits are exactly those
two sets. The last marked generator belongs to every
one of these omissions and therefore preserves both
sides of every tree cut. Tree cuts distinguish every
pair of vertices, so it fixes every vertex. This
contradicts its nonidentity. When $d=1$, degree one
likewise makes the sole mark trivial. Therefore the
degree is at least $d+1=x-4$, as claimed.
\end{proof}

\begin{theorem}
With the actual transitive witness and transport above,
the last active label $x$ satisfies
\[
 x<r-1,\qquad |\operatorname{Fix}(G_{>x})|\ge2x+1.
\]
\end{theorem}
\begin{proof}
The eight points $D\cup\{\beta,z_0\}$ and the two
nontrivial $K$-orbits $O,T$ are disjoint. They are
all fixed by every mark after $x$. The preceding
lemma gives at least
\[
                         8+2(x-4)=2x              \tag{1}
\]
fixed points.

Suppose first that $x<r-1$ but the asserted stronger
bound fails. Then the fixed set $F$ consists exactly
of these eight points and the two $K$-orbits, with
$|O|=|T|=x-4$. The involution $\delta$ centralizes
$K$ and $G_{>x}$, and fixes the eight points and
$T$ pointwise. It preserves $F$, so it must preserve
$O$ as well. Its restriction to $O$ is the
restriction of the actual $\tau\in G_j$. It
commutes with $\rho_j$ and preserves the unique
ambient $j$-edge there. Since its endpoints belong
to different $G_j$-orbits, each endpoint is fixed.
Centralization of the transitive $K$-action now
makes $\delta$ fix all of $O$.

Thus $\delta$ fixes $F$ pointwise. The set $F$
is $K$-invariant by its displayed decomposition.
It is preserved by $\rho_5$, which commutes with
$G_{>x}$ because $x\ge6$, and by $\sigma$, which
also centralizes that suffix. The first three
marks preserve $D$, and $\tau=(6\ \beta)\delta$
preserves $F$. Hence every generator of $G$
preserves $F$. Transitivity would force $F=\Omega$,
but then the odd involution $\delta$ would be
identity. This contradiction proves the stronger
bound whenever $x$ is nonterminal.

It remains to exclude $x=r-1$. Write $N=|O|=|T|$.
The head and these two orbits occupy at least
$8+2(r-5)=2r-2$ points, leaving at most three.
If $\delta$ preserves $O$, the same unique-edge
argument fixes it pointwise. The odd involution
$\delta$ would then be supported on at most three
points, hence be a literal transposition. This is
excluded by the complete acquired exterior-pair
obstruction.

Otherwise $\delta$ sends $O$ to a distinct orbit
$V$ of $H$, because it centralizes $H$. It is
also an actual $K$-orbit, equivariantly isomorphic
to $O$. It is disjoint from $T$, which $\delta$
fixes pointwise, and from the eight fixed points.
Consequently $N\le3$. Since $N\ge r-5$ and
$x=r-1\ge6$, only the following cases can occur.
For $r=7$, necessarily $N=2$, since three copies
of degree three would not fit. The head and the
three copies occupy fourteen of the fifteen points.
The involution $\delta$ exchanges $O,V$ by two
transpositions, fixes $T$ and the eight head points,
and fixes the one remaining point. It is even,
a contradiction.

For $r=8$, necessarily $N=3$, and the eight head
points and the three copies exhaust all seventeen
points. On each copy the two active involutions
$\rho_6,\rho_7$ are distinct transpositions generating
natural $S_3$. Their equivariant identifications
make $K$ diagonal on three natural rows. The
centralizer of natural $S_3$ on one row is trivial,
and there is a unique equivariant bijection between
any two rows. Thus every permutation centralizing
$K$ acts on the row union solely by permuting the
three rows.

The row action of $\delta$ exchanges $O,V$ and
fixes $T$. The row action of $\sigma$ must exchange
$T$ with one of $O,V$. Otherwise $O\cup V$ would
be invariant under every generator: $H$ preserves
its two actual orbits, $\tau$ exchanges them through
$\delta$, $\sigma$ preserves their union, and the
first three marks fix it. This contradicts
transitivity. Thus the two row involutions are
distinct transpositions. Their product is a
three-cycle on the rows.

On the fixed eight-point set, $\tau$ and $\sigma$
are respectively $(6\ \beta)$ and $(\beta\ z_0)$,
with the other five points fixed. Their product
is likewise a three-cycle. Hence the dihedral
group $\gen{\tau,\sigma}$ has its nonidentity
rotations nontrivial both on the three-point
core and on the row union. Its reflections have
nontrivial row transpositions. It cannot contain
a pure three-cycle on $\{6,\beta,z_0\}$, contrary
to the acquired actual overlap theorem. This
excludes the final terminal case and proves
the theorem.
\end{proof}

\paragraph{Attribution and boundary.}
The strong witness colour, whole-orbit transport,
actual head anchor and pure core cycle are acquired.
The commuting-factor proof, omission forest and
complete fixed-set argument supply the new separator.
The terminal three-row argument reuses the acquired
odd-degree row mechanism with its actual hypotheses
verified here; no earlier row-assignment theorem
is invoked with changed assumptions. The second
witness need not be locally unique in its restricted
$K$-action, and $\rho_5$ is not assumed to fix it.
Full global source routing and attachments remain.
No whole-original, question, numerical or formal
credit is claimed.
\end{document}
